% --- Extended Chapters ---
\newpage
\begin{center}
{\fontsize{14}{17}\bfseries Chapter 7\\[6pt]}
{\fontsize{16}{20}\bfseries SYSTEM IMPLEMENTATION DETAILS\\[10pt]}
\end{center}
\addcontentsline{toc}{section}{Chapter 7: System Implementation Details}

\subsection*{7.1 Data Pipeline for E-Commerce Products}
The data pipeline is the foundational layer of the ConvoShop system, responsible for acquiring, cleaning, normalizing, and standardizing product data from disparate sources. E-commerce datasets, particularly those sourced from public repositories like Kaggle and Hugging Face, are often noisy, inconsistent, and contain numerous missing values. 

\subsubsection*{7.1.1 Mobile Phones Dataset Processing}
The mobile phones dataset, sourced from Abdul Malik's 2025 Mobiles Dataset on Kaggle and augmented with Nadirova's Phone\_SpecsDataset\_25K on Hugging Face, initially contained thousands of records with varying degrees of completeness. The processing pipeline involved several critical stages. First, data ingestion standardized the column names and data types, ensuring that numeric fields like battery capacity (mAh), RAM (GB), and storage (GB) were correctly parsed as integers rather than strings. 

Second, the cleaning phase addressed missing values and structural inconsistencies. A common issue was the inclusion of tablet records within the phone dataset; these were identified and filtered out based on screen size thresholds and categorical tags. 

Third, the merging phase aligned products across the two datasets. Given that the same phone model (e.g., "Samsung Galaxy S24 Ultra") might be named slightly differently in each source, fuzzy string matching and brand-model normalization heuristics were employed to deduplicate records and enrich the final dataset with the most accurate specifications and image references available. The final cleaned dataset for mobile phones comprised 271 high-quality, verified records.

\subsubsection*{7.1.2 Laptops Dataset Processing}
The laptop dataset pipeline followed a similar architectural pattern but introduced category-specific challenges. A primary challenge was the currency normalization. The raw dataset contained prices in Euros (EUR), whereas the system's operational currency for the Indian market is Indian Rupees (INR). A fixed-rate conversion mechanism was implemented during the cleaning phase to standardize all price fields to INR, enabling accurate cross-category budget filtering. 

Furthermore, laptop titles are notoriously noisy, often including the entire specification list (e.g., "Dell XPS 13 9310 Core i7 16GB 512GB SSD"). The data engineering process involved extracting structured features—such as processor generation, RAM capacity, and storage type—from these noisy titles using regular expressions, resulting in a finalized catalog of 1,453 laptop records.

\subsubsection*{7.1.3 Smartwatches Dataset Processing}
The smartwatch dataset required careful handling of brand normalization and estimated values. Given the fragmented nature of smartwatch specifications across different manufacturers, many features were either missing or inconsistently reported. The pipeline distinguished between verified original values and estimated values (e.g., inferring battery life from battery capacity and display type). This distinction is crucial for maintaining transparency in the recommendation engine, ensuring that the LLM does not confidently assert estimated specifications as facts. The smartwatch pipeline yielded 830 finalized records.

\subsection*{7.2 Unified Product Knowledge Base}
The culmination of the category-specific pipelines is the Unified Product Knowledge Base. This centralized repository combines the 271 phone, 1,453 laptop, and 830 smartwatch records into a standardized schema. 

For each product, a comprehensive natural-language description was procedurally generated by concatenating key specifications and features into a coherent paragraph. These descriptions serve as the foundational text for the semantic retrieval process. 

\subsection*{7.3 Frontend Architecture: React.js and Vite}
The conversational user interface is built using React 19, initialized with Vite 8 for rapid development and optimized production builds. The frontend architecture is entirely component-based, leveraging functional components and React Hooks (useState, useEffect, useMemo, useRef) for state management and side-effect handling.

\subsubsection*{7.3.1 Application State Management}
Local state management is utilized for handling the conversational session context, user inputs, and streaming responses from the backend. The chat interface dynamically renders user queries and system responses, updating incrementally as tokens are received via Server-Sent Events (SSE) from the FastAPI backend.

\subsubsection*{7.3.2 Data Parsing and Local Storage}
For demonstration purposes and rapid prototyping, the frontend utilizes PapaParse 5 to read an Amazon products CSV (containing approximately 42,000 rows) directly in the browser. User sessions, including demo logins, saved preferences, and virtual shopping carts, are persisted using the browser's native localStorage and sessionStorage APIs. 

\subsubsection*{7.3.3 Styling and Aesthetics}
The aesthetic design of ConvoShop relies on plain CSS (App.css, index.css) without the overhead of external UI component libraries like Tailwind CSS or Bootstrap. Typography is carefully selected, utilizing Google Fonts' Bricolage Grotesque for headings and Instrument Sans for body text, contributing to a clean, modern, and professional e-commerce advisory interface.

\subsection*{7.4 Backend Architecture: FastAPI and Python}
The backend infrastructure is developed using Python, orchestrated by the FastAPI framework. FastAPI was selected for its high performance, native support for asynchronous programming (asyncio), and automatic generation of interactive API documentation (Swagger UI).

\subsubsection*{7.4.1 AI Pipeline and Embedding Generation}
The AI pipeline heavily relies on the sentence-transformers library to generate dense vector embeddings from the natural-language product descriptions. The embedding generation process is executed asynchronously, converting user queries into high-dimensional vectors that capture the semantic intent behind the natural language. 

\subsubsection*{7.4.2 Vector Retrieval Integration}
For the vector retrieval component, the system utilizes FAISS (Facebook AI Similarity Search) and ChromaDB to perform highly efficient Approximate Nearest Neighbor (ANN) searches over the product knowledge base. The integration allows the system to scale to millions of product embeddings while maintaining sub-100 millisecond retrieval latency.

\newpage
\begin{center}
{\fontsize{14}{17}\bfseries Chapter 8\\[6pt]}
{\fontsize{16}{20}\bfseries SYSTEM EVALUATION AND TESTING\\[10pt]}
\end{center}
\addcontentsline{toc}{section}{Chapter 8: System Evaluation and Testing}

\subsection*{8.1 Evaluation Methodology}
Evaluating a Retrieval-Augmented Generation (RAG) system presents unique challenges compared to traditional recommendation engines. The evaluation must assess both the accuracy of the retrieval component (the Vector Database) and the quality of the generated response (the LLM). For ConvoShop, a 70-question evaluation dataset was developed, encompassing a wide range of user intents, from specific technical requirements to broad budget constraints.

\subsection*{8.2 Retrieval Metrics}
The retrieval performance was evaluated using standard information retrieval metrics:
\begin{itemize}
    \item \textbf{Precision@K:} Measures the proportion of relevant products in the top-K retrieved results. ConvoShop demonstrated high Precision@10, successfully retrieving products that semantically matched the complex, multi-constraint queries.
    \item \textbf{Recall@K:} Evaluates the system's ability to find all relevant products within the database. The integration of sentence-transformers and FAISS significantly improved Recall compared to keyword-based baselines, particularly for queries employing synonymous or descriptive language (e.g., "good for photography" retrieving phones with high megapixel counts and advanced sensor sizes).
\end{itemize}

\subsection*{8.3 Constraint Accuracy}
A critical requirement for e-commerce advisory is strict adherence to explicit user constraints, particularly budget and brand filters. The evaluation assessed constraint accuracy by running tests with strict maximum/minimum specification queries (e.g., "strictly under 30,000 INR"). The RAG architecture, combined with explicit metadata filtering prior to the LLM generation phase, ensured that recommended products strictly adhered to numerical constraints, eliminating a common failure mode of standalone LLMs.

\subsection*{8.4 Hallucination Mitigation}
A primary objective of ConvoShop is the reduction of LLM hallucination—the generation of plausible but factually incorrect product information. By strictly grounding the Gemini 1.5 Flash model in the retrieved product metadata (the context document), hallucination rates were dramatically reduced. Spot-checks of the generated responses against reliable manufacturer sources confirmed that specifications, prices, and features presented to the user were factually accurate and directly traceable to the underlying database.

\subsection*{8.5 Performance and Scalability}
Latency analysis confirmed that the FastAPI backend, coupled with the FAISS/ChromaDB vector search, consistently delivered retrieval results within the specified 100ms threshold under normal load conditions. The asynchronous design of the backend ensures that concurrent user requests are handled efficiently, preventing blocking operations during the LLM inference phase. The system architecture is inherently scalable, capable of accommodating an expanding catalog of products across new categories (e.g., home appliances, apparel) with minimal architectural modifications.

